% Standalone problem document, generated from the adjacent README.md.
% Compile with XeLaTeX. Mathematical content is maintained in Markdown.
\documentclass[11pt,a4paper]{article}
\usepackage[fontsize=10.5pt]{scrextend}
\usepackage[margin=22mm,headheight=15pt,headsep=9mm,footskip=11mm]{geometry}
\usepackage{fontspec}
\setmainfont{texgyrepagella}[Extension=.otf,UprightFont=*-regular,BoldFont=*-bold,ItalicFont=*-italic,BoldItalicFont=*-bolditalic]
\setsansfont{texgyreheros}[Extension=.otf,UprightFont=*-regular,BoldFont=*-bold,ItalicFont=*-italic,BoldItalicFont=*-bolditalic]
\setmonofont{texgyrecursor}[Extension=.otf,UprightFont=*-regular,BoldFont=*-bold,ItalicFont=*-italic,BoldItalicFont=*-bolditalic,Scale=MatchLowercase]
\usepackage{amsmath,amssymb}
\usepackage{unicode-math}
\setmathfont{texgyrepagella-math.otf}
\usepackage{xcolor,microtype,enumitem,needspace,fancyhdr}
\usepackage{bookmark}
\definecolor{ink}{HTML}{17344B}
\definecolor{accent}{HTML}{197D80}
\definecolor{muted}{HTML}{596773}
\hypersetup{colorlinks=true,linkcolor=accent,urlcolor=accent,pdftitle={TR-16 - Monotonicity
of the average number of critical rank-one
approximations},pdfauthor={Open Problems in Numerical Linear Algebra}}
\urlstyle{same}
\setlength{\parindent}{0pt}
\setlength{\parskip}{5pt plus 1pt minus 1pt}
\linespread{1.04}
\setlength{\emergencystretch}{3em}
\widowpenalty=10000
\clubpenalty=10000
\displaywidowpenalty=10000
\allowdisplaybreaks[1]
\setlist{nosep,leftmargin=1.5em}
\providecommand{\tightlist}{\setlength{\itemsep}{0pt}\setlength{\parskip}{0pt}}
\providecommand{\pandocbounded}[1]{#1}
\pagestyle{fancy}
\fancyhf{}
\fancyhead[L]{\sffamily\footnotesize\color{muted}OPEN PROBLEMS IN NUMERICAL LINEAR ALGEBRA}
\fancyhead[R]{\sffamily\bfseries\footnotesize\color{accent}TR-16}
\fancyfoot[L]{\parbox[b]{0.9\textwidth}{\sffamily\fontsize{8}{10}\selectfont\color{muted}Editorial ratings; a bounded search cannot certify that a problem remains unsolved.\\Literature check: 2026-09-10}}
\fancyfoot[R]{\sffamily\footnotesize\color{muted}\thepage}
\renewcommand{\headrulewidth}{0.3pt}
\renewcommand{\headrule}{\hbox to\headwidth{\color{accent}\leaders\hrule height \headrulewidth\hfill}}
\makeatletter
\renewcommand{\section}{\Needspace{5\baselineskip}\@startsection{section}{1}{0pt}{12pt plus 2pt minus 2pt}{4pt}{\sffamily\bfseries\large\color{ink}}}
\renewcommand{\subsection}{\Needspace{4\baselineskip}\@startsection{subsection}{2}{0pt}{9pt plus 1pt minus 1pt}{3pt}{\sffamily\bfseries\normalsize\color{ink}}}
\makeatother
\setcounter{secnumdepth}{0}
\begin{document}
{\sffamily\small\color{muted}Tensor computations\par}
\vspace{3pt}
{\raggedright\hyphenpenalty=10000\exhyphenpenalty=10000\sffamily\bfseries\fontsize{21}{24}\selectfont\color{ink}Monotonicity
of the average number of critical rank-one approximations\par}
\vspace{7pt}
{\sffamily\small\textbf{Difficulty:} challenging\quad\textbf{Importance:} interesting
to specialist\par}
{\sffamily\small\bfseries\color{accent}Status: Open\par}
\vspace{4pt}
{\color{accent}\hrule height 0.8pt}
\vspace{6pt}
\textbf{Rating rationale:} Challenging because a dimension comparison of
Gaussian critical-point expectations needs control beyond the known
integral formula; specialist importance concerns a particular statistic
of rank-one approximation landscapes.

\subsection{Statement}\label{statement}

For integers \(p\ge3\) and \(2\le n_1\le\cdots\le n_p\), let
\(G\in\bigotimes_{i=1}^p\mathbb R^{n_i}\) have independent standard
normal entries. Let \(X\) be the smooth manifold of nonzero real
rank-one tensors in this space, and let \(N(G)\) be the number of
critical points on \(X\) of \[
Z\longmapsto\|G-Z\|_F^2.
\] Each critical tensor is counted once, regardless of how its vector
factors are scaled. Define \(a(n_1,\ldots,n_p)=\mathbb E[N(G)]\); the
critical-point count is finite almost surely.

If \[
n_p-1>\sum_{i=1}^{p-1}(n_i-1),
\] is it always true that \[
a(n_1,\ldots,n_{p-1},n_p)
\le a(n_1,\ldots,n_{p-1},n_p-1)?
\] The expectation on the right uses a standard Gaussian tensor in the
smaller space.

\subsection{Relevance}\label{relevance}

Rank-one approximation algorithms navigate a landscape of stationary
points. This conjecture predicts how the average number of such points
changes when one mode becomes much larger than the others.

\subsection{References}\label{references}

\begin{enumerate}
\def\labelenumi{\arabic{enumi}.}
\tightlist
\item
  J. Draisma and E. Horobeţ, \emph{The average number of critical
  rank-one approximations to a tensor}, Linear Multilinear Algebra 64
  (2016), 2498--2518.
  \href{https://doi.org/10.1080/03081087.2016.1164660}{DOI};
  \href{https://arxiv.org/pdf/1408.3507}{primary preprint}, Conjecture
  1.3, p.3; §1 defines the Gaussian average and Theorem 1.1 gives an
  integral formula.
\item
  S. Friedland and G. Ottaviani, \emph{The Number of Singular Vector
  Tuples and Uniqueness of Best Rank-One Approximation of Tensors},
  Found. Comput. Math. 14 (2014), 1209--1242.
  \href{https://doi.org/10.1007/s10208-014-9194-z}{DOI};
  \href{https://arxiv.org/pdf/1210.8316}{primary preprint}, Theorem 1,
  for the corresponding complex critical-point count.
\end{enumerate}

\subsection{Status check --- 2026-09-10}\label{status-check-2026-09-10}

Rechecked \href{https://arxiv.org/pdf/1408.3507}{Draisma--Horobeţ,
Conjecture 1.3}, including the direction of the inequality, and searched
by the authors, critical-rank-one counts and monotonicity. No proof or
counterexample to the real Gaussian statement was located. The matrix
case is excluded here, and stabilization of complex critical-point
counts does not settle this expectation. No recent explicit primary
reaffirmation was found.
\end{document}
